\documentclass[11pt]{article}

\usepackage[a4paper,margin=1in]{geometry}
\usepackage{amsmath,amssymb,mathtools}
\usepackage{booktabs}
\usepackage{xcolor}
\usepackage{enumitem}
\usepackage{titlesec}
\usepackage{hyperref}

\definecolor{ink}{HTML}{16181D}
\definecolor{accent}{HTML}{1F6F5C}
\definecolor{accent2}{HTML}{7A3B69}
\definecolor{rule}{HTML}{C9CDD4}

\hypersetup{colorlinks=true, linkcolor=accent, urlcolor=accent, citecolor=accent}

\titleformat{\section}{\color{ink}\Large\bfseries}{\thesection}{0.6em}{}
\titleformat{\subsection}{\color{ink}\large\bfseries}{\thesubsection}{0.6em}{}
\titlespacing*{\section}{0pt}{1.4em}{0.6em}

\newcommand{\dt}{\Delta t}
\newcommand{\rate}{\dot{\gamma}}
\newcommand{\norm}[1]{\left\lVert #1 \right\rVert}
\newcommand{\note}[1]{\vspace{0.3em}\noindent{\color{accent2}\small\textsf{#1}}\vspace{0.3em}}

\title{\vspace{-2em}\textbf{NNUI\_CELL}\\[0.3em]
\large Addendum v0.3 --- Workstream-A Realizations and the\\
\large Discrete Interface-Tension Design\\[0.6em]
\normalsize Supplement to the v0.2 Specification (which remains canonical)}

\author{R.~Forrestall (``Heat'')\\ \normalsize MIVIM Development --- \textit{Standards in Novel Design}}
\date{July 2026}

\begin{document}
\maketitle

\begin{abstract}
\noindent
This addendum records what building workstream~A --- the agent-facing API, its
integration demo, and the tessellation stage --- forced into the discrete
realization of NNUI\_CELL. The continuum theory of v0.1/v0.2 is again
unchanged. Five corrections (C5--C9) join the normative set: the discrete
realization of $E_{\text{fill}}$ with a wall-coverage stall predicate, plastic
post-yield rest adaptation, bulk-transport binding, the latch's exact
$\xi{=}1$ assignment, and radius conservation under lattice dilation. The
final section is a \emph{design}, not a result: the discrete realization of
$E_{\text{cohesion}}$ as boundary curvature flow, motivated by three measured
failures of naive compactness mechanisms, with binary validation exits (T1--T4)
that implementation must pass before the design is admitted as normative.
Evidence levels are marked throughout: C5--C9 are validated
(suite-green, screen-verified); the interface-tension design is internal
reasoning awaiting its experiment.
\end{abstract}

\noindent\textcolor{rule}{\rule{\linewidth}{0.8pt}}

\section{Position of this document}

The v0.2 specification remains the canonical statement of the model. Everything
here concerns \emph{discrete realization} --- the same relationship v0.2's
corrections C1--C4 bear to v0.1. Workstream~A validated the API surface
(milestones M1--M4: atomic grammar submission, the event-driven read path,
LLM-driven layout, portable geometry-free cell records) and then confronted the
layout semantics the prototype had deferred: the medium must actually fill its
vessel, sizes must actually be solved from material state, and bound groups
must actually cohere. Realizing those claims produced C5--C9 and exposed the
one energy term still without a discrete form, $E_{\text{cohesion}}$.

\section{Corrections C5--C9 (validated, normative)}

\begin{enumerate}[leftmargin=2.2em,label=\textbf{C\arabic*.},start=5]

\item \textbf{$E_{\text{fill}}$ is a quasi-static drive with a wall-coverage
stall.} The fill phase advances matter outward at a rate that yields against
solidity, $g = g_0\,(1-S)$, entering the predictor as displacement (per C2).
While filling, the frame is effectively accelerating (the boot pour): $\xi$
rebuild is suspended per the \S7.2 boundary condition, so the medium flows
freely into void. Termination is \emph{pressure stall}, not an area target ---
and the stall predicate may not fire while boundary spans remain unclaimed:
edge-to-edge is part of $E_{\text{fill}}$'s definition
($\rho\!\to\!1$ over $\Omega$), not a preference. Measured: with an area-target
stall the settled screen claimed 50--75\% of the vessel; with the wall-coverage
predicate, 95--100\%. Near walls the drive direction blends toward the local
void (matter creeps along walls; corner zones are owned by void-seeking);
applied globally, void-seeking pulls cells into elongated shards --- measured,
rejected, and the restriction is therefore normative.

\item \textbf{Post-yield flow is plastic.} While a layout region is awake
(yielded), its material rest lengths follow realized geometry at a rate that
yields against solidity, $\dot r_{\text{rest}} \propto (1-S)$; at rest the
rests hold, elastic below yield. This is the discrete realization of
shear-thinning post-yield flow ($n{<}1$): committed matter gets out of the way
and takes its new shape; stiff matter keeps its shape. Without plasticity the
fill phase fights permanent elastic memory of the spawn lattice and the medium
cannot tile.

\item \textbf{Binds transport bulk.} The deliberate \texttt{bind} operation is
interfacial: it must move \emph{regions}, not points. A single-anchor spring
fails in a plastic medium --- the anchor migrates through the region while the
bulk never moves (measured: bound chains left 170--240\,px gaps). A
stiffness-scaled attraction fails at rest --- compliance stiffening turns it
into a slingshot (the compliance-at-rest failure class, third sighting;
measured eruption at 140\,u/s). The normative form: a bounded, constant-scale
quasi-static drive on \emph{all} particles of both bound regions, active until
the pair sits adjacent, with stiff contacts preventing overlap. Cohesive
groups then pack, solved.

\item \textbf{The latch assigns $\xi = 1$ exactly.} \S7.3 writes ``set
$\xi_i = 1$'' and means it. A latch that merely stops motion lets each
region's $\xi$ carry its private settle history --- measured as
identically-declared cells exporting different material state in their
portable records (settle history leaking into the wire format). On latch
engage, $\xi \gets 1$: structure is fully rebuilt at rest, and exported
material state is a function of declaration and physics, not of arrival
order.

\item \textbf{Radius conservation under dilation.} A dilated lattice is not a
porous point-cloud: when lattice spacing grows, effective particle radius
grows with it, $r_p = \tfrac{1}{2}\,\bar\ell$ (with $\bar\ell$ the region's
mean lattice spacing), clamped to a sane range. Without this, regions finger
through each other's interiors and settle interleaved (measured: hull-overlap
``stacked sheets''). Two consequences are part of the correction: the C1
topological-guard tolerance must be \emph{proportional} to the contact rest
(an absolute tolerance misclassifies ordinary residual prestress at dilated
scale), and contact-search structures must follow the radius. The material
radius is exposed through the read path so a renderer draws true material
extent --- physical seams are seams, not voids.

\end{enumerate}

\note{Status: C5--C9 are validated --- the engine suite (35+ tests including
all seven invariants, the four milestone harnesses, frame budget, and
API-boundary determinism) is green with all five active, and the settled
screen is verified programmatically at 95--100\% vessel claim.}

\section{Design: $E_{\text{cohesion}}$ as boundary curvature flow}

\subsection{The residual defect and why it points here}

With C5--C9 active, one defect class remains, measured and reproducible:
\emph{interleaving}. Arms of soft regions, stretched during early growth,
freeze across neighbouring cells; per-region mean radius cannot seal
non-uniformly stretched arms; vessel corners adjacent to such arms stay
partially unclaimed. Nothing in the realized model penalizes elongation. In
the continuum theory that penalty exists and has a name:
\[
E_{\text{cohesion}} \;=\; w_c \sum_{(i,j)\in\partial} \sigma_{ij}\,\ell_{ij},
\]
interface \emph{length} minimization --- surface tension. Foam tessellates
because pressure fills while tension keeps every cell compact. We have built
the pressure and not the tension.

\subsection{Failed candidates (measured, excluded)}

\begin{itemize}[leftmargin=1.4em]
\item \emph{Global void-seeking} as a shape force: produces taffy shards ---
it is a fill term, not a tension term, and only survives restricted to the
wall band (C5).
\item \emph{Radial roundness tension} (pull each particle toward the region's
mean radius): isotropic, so it fights wall flattening and the fill drive;
measured to collapse vessel claim to 70\% and break solved sizes. The error is
structural: it penalizes \emph{radius variance}, but a wall-flattened tile has
high radius variance and \emph{low perimeter}. The energy says perimeter.
\end{itemize}

\subsection{The design}

Realize interface tension as \textbf{discrete curvature flow on region
boundaries}:

\begin{enumerate}[leftmargin=1.6em]
\item \textbf{Boundary detection.} Per frame, a particle is a boundary
particle of its region iff it has fewer than $B$ same-region neighbours
within radius $\rho_n$ (existing spatial grid; index-order scan;
deterministic).
\item \textbf{Discrete curvature.} For boundary particle $i$, let $m_i$ be
the mean position of its same-region \emph{boundary} neighbours within
$\rho_b$. The discrete curvature vector is
\[
\kappa_i \;=\; m_i - x_i ,
\]
the boundary Laplacian: on a straight (wall-flattened or shared-interface)
edge $\kappa_i \approx 0$; on a convex bulge $\kappa_i$ points inward; in a
concave dent it points outward.
\item \textbf{Tension step.} Quasi-static displacement (C2), boundary
particles of awake regions only:
\[
\Delta x_i \;=\; \min\!\big(w_c'\,\norm{\kappa_i},\; c_{\max}\big)\,
\hat\kappa_i\,\dt .
\]
Interior particles are untouched; the arrest phase and honest strain
measurement are untouched; contacts and walls still project afterwards.
\end{enumerate}

\subsection{Why this succeeds where the candidates failed}

Curvature flow minimizes perimeter directly: flat interfaces (shared seams,
wall contact) carry no force, so it cannot fight wall conformity or tiling;
bulges retract; dents fill; and a thin shard arm --- whose boundary curvature
is everywhere high and inward --- retracts into the body of its cell. The
predicted equilibrium is exactly the paper's stated consequence of
$E_{\text{cohesion}}$: compact polygonal tiles, flat shared interfaces, and
rounded free corners --- coalescence and menisci. Because the force vanishes
on flat boundaries it terminates naturally under the deadband and latch; no
new stall machinery is required. It is diffusive and therefore stable at a
bounded step ($c_{\max}$), and deterministic by construction.

\subsection{Validation exits (binary, pre-registered)}

Implementation is admitted as normative only if all four pass:

\begin{center}
\begin{tabular}{@{}p{0.07\linewidth}p{0.85\linewidth}@{}}
\toprule
\textbf{T1} & \emph{No interleaving.} For every reference layout, pairwise
convex-hull intersection area between settled cells is $<2\%$ of the smaller
cell's area. \\
\textbf{T2} & \emph{Corners claimed.} All four vessel corners pass the
50\,px / 30\% matter standard in every layout. \\
\textbf{T3} & \emph{Nothing regresses.} Vessel claim stays $\geq 90\%$ lit;
solved-size ratios, bind adjacency, phenotype expression, frozen-raster rest,
and the \S9 resize demonstration all hold. \\
\textbf{T4} & \emph{The suite is the suite.} All engine tests green;
bit-identical determinism preserved; the flicker-free sustained-rest window
holds. \\
\bottomrule
\end{tabular}
\end{center}

\note{Status: design only --- internal reasoning strength. T1--T4 are the
experiment; the design is wrong if they fail, and the failure record joins
this document's lineage the same way C1--C9 earned theirs.}

\section{Design revision D2: constraint-based tessellation}

\note{Added after implementation of the curvature-flow design. Status of the
first pass, honestly: curvature flow itself behaved as designed (arms
retract, compact cells) and survives; but the accompanying drives ---
fill, sibling exclusion, bind transport --- were bounded competing
displacements, and five exclusion formulations at multiple authorities all
equilibrated \emph{wherever their caps happened to balance} (pairwise hull
overlap hovering near 4\% against T1's 2\% bar, with any strengthening
breaking bind adjacency, solved sizes, or fill). The measured conclusion:
\textbf{bounded competing drives cannot express the tessellation
equilibrium.} Two genuine corrections were also forced and are kept: the
\emph{phase gate} (the fill phase is supra-yield by definition --- no
deadband arrest, no latching mid-pour; without it a stiff cell's sub-yield
growth is reverted every frame and it latches at spawn size), and the
\emph{wall-wetting exemption} (tension acts on free surfaces; matter
pressed against the vessel is wetting it, and tension at a corner tip
evacuates the wedge $E_{\text{fill}}$ claims).}

\subsection{The reformulation}

Move the three competing terms \emph{into} the XPBD projection loop, where
equilibria are solved rather than raced:

\begin{enumerate}[leftmargin=1.6em]
\item \textbf{Fill as a per-region area constraint} (the paper's
volume/occupancy constraint, \S5.1). Per region $r$:
$C_r = A_r(\mathbf{x}) - A_r^{*}$, with $A_r$ the area of the region's
ordered boundary polygon (convex hull at this fidelity) and targets solved
from material state, $A_r^{*} \propto (1-S_r)^2$ over the vessel area.
The gradient on hull vertex $v_i$ is the standard polygon-area gradient
$\tfrac{1}{2}\,\mathrm{perp}(v_{i+1}-v_{i-1})$; one scalar multiplier per
region per frame. Because the multiplier accumulates inside the solve,
elastic pullback cannot starve growth --- the stuck-stiff-cell failure
class dies structurally. Sizes remain \emph{solved}: authority lives in
the targets, not in growth-rate races.
\item \textbf{Exclusion largely for free.} Under hull-area constraints,
wrapping is self-defeating: enclosed foreign space counts against the
wrapper's own area budget while the enclosed cell's constraint pushes back
from inside. Particle non-penetration (C1) plus conserved radii (C9)
handle contact; an explicit exclusion survives only as a light correction
for measured bites, if any remain.
\item \textbf{Bind as a plain XPBD distance constraint} (few anchor pairs,
constant soft compliance --- never stiffness-scaled, per C7). With area
protected by its own constraint, a region can no longer comply by
deforming around a migrating anchor: translation is the only compliant
path, so binds transport bulk \emph{because of} the constraint coupling,
not despite it.
\end{enumerate}

Curvature flow (\S3), the phase gate, the wall-coverage stall (C5), plastic
rests (C6), and radius conservation (C9) are unchanged. The same
pre-registered exits T1--T4 gate this revision; T1 remains recorded-open
until it passes with full fill.

\section{New constants and dials}

\begin{center}\small
\begin{tabular}{@{}lll@{}}
\toprule
Symbol & Meaning & Status \\
\midrule
$g_0$ & fill drive scale (yields against $S$) & tuned (C5) \\
$w_{\text{band}}$ & wall band for void-seeking & tuned (C5) \\
$\theta_{\text{wall}}$ & wall-coverage stall threshold ($>90\%$) & tuned (C5) \\
$p$ & plastic rest-adaptation rate ($\propto 1-S$) & tuned (C6) \\
$g_b$ & bind bulk-drive scale (bounded, constant) & tuned (C7) \\
$r_p$ & effective particle radius, $\tfrac{1}{2}\bar\ell$ & derived (C9) \\
$B, \rho_n, \rho_b$ & boundary detection neighbourhood & design (T1--T4) \\
$w_c', c_{\max}$ & curvature-flow weight and step clamp & design (T1--T4) \\
\bottomrule
\end{tabular}
\end{center}

\noindent As with the v0.2 operating point, the corrections are qualitative
and forced by measured failures; the constants are tuned, not derived.

\vspace{1em}
\noindent\textcolor{rule}{\rule{\linewidth}{0.8pt}}
\vspace{0.4em}

\noindent{\small \textit{MIVIM Development --- Standards in Novel Design.}
Addendum v0.3, July 2026; supplements the v0.2 specification, which remains
canonical. All results referenced herein are reproducible from the repository
sources (\texttt{engine/}, \texttt{demo/}) and their automated gates.}

\end{document}
